\chapter{Introduction à la Gestion de Crise}
\label{chap:introcrise}

La gestion de crise est l'ultime rempart. Lorsque la prévention a échoué et que l'incident dépasse les capacités de réponse opérationnelle, l'entreprise entre en zone de turbulence. C'est un moment de vérité pour la gouvernance.

\section{Présentation : De l'Incident à la Crise}

\subsection{La Rupture de l'Équilibre}

Une panne informatique ou une attaque peut être un incident gérable. Mais elle devient une \textbf{crise} lorsqu'elle menace la pérennité de l'organisation, sa réputation, ou met en danger des vies humaines.

\begin{boxdef}
\textbf{Crise :} Événement soudain, inattendu et menaçant qui perturbe gravement le fonctionnement normal de l'organisation et nécessite des décisions rapides sous forte pression.
\end{boxdef}

La crise se caractérise souvent par :
\begin{itemize}
    \item \textbf{L'incertitude :} On ne maîtrise pas toute l'information.
    \item \textbf{L'urgence :} Le temps manque pour décider.
    \item \textbf{La pression médiatique :} L'extérieur (presse, réseaux sociaux) regarde et juge.
\end{itemize}

\subsection{Distinction Incident vs Crise}

Il est fondamental de ne pas confondre "Incident Majeur" et "Crise".

\begin{table}[htbp]
\centering
\begin{tabular}{p{4cm}p{4cm}p{4cm}}
    \toprule
    \textbf{Critère} & \textbf{Incident Majeur} & \textbf{Crise} \\
    \midrule
    \textbf{Focus} & Technique & Stratégique / Image \\
    \textbf{Gestion} & DSI / Technique & Direction Générale \\
    \textbf{Communication} & Technique (Interne) & Médiatique / Externe \\
    \textbf{Durée} & Heures à Jours & Semaines à Mois \\
    \textbf{Impact} & Arrêt activité & Survie de l'entreprise \\
    \bottomrule
\end{tabular}
\caption{Comparaison Incident Majeur et Crise}
\end{table}

\begin{boxlizbiz}
\textbf{Le "Big Bang" chez LiZbiZ :}
Jusqu'ici, le Ransomware était un incident technique (Module 3).
Ce matin, un journaliste du "Monde" appelle la direction : "Nous avons en notre possession des preuves que les données de vos clients sont en vente sur le Dark Web. Que faites-vous ?".
\textbf{Bascule :} L'incident technique devient une \textbf{Crise Médiatique et Juridique}.
\end{boxlizbiz}

% ------------------------------------------------------------------------------
\section{Concepts Clés}
% ------------------------------------------------------------------------------

\subsection{Le Seuil de Rupture}

Toute entreprise a une capacité d'absorption des chocs. La crise survient quand ce seuil est dépassé.
Ce seuil dépend de la culture de l'entreprise, de sa trésorerie, et de sa réputation préalable.

\subsection{Les Impacts Multidimensionnels}

Une crise cyber n'est jamais uniquement cyber. Elle a des répercussions en cascade :

\begin{enumerate}
    \item \textbf{Impact Financier :} Perte de CA, coût de la réparation, amendes (CNIL), baisse du cours de bourse.
    \item \textbf{Impact Réputationnel :} Perte de confiance des clients, dégradations de la marque employeur.
    \item \textbf{Impact Juridique :} Mise en cause de la responsabilité des dirigeants (infractions pénales possibles).
    \item \textbf{Impact Humain :} Stress intense, burn-out des équipes techniques, conflits internes.
\end{enumerate}

\subsection{L'Instinct de Survie}

En temps de crise, les processus habituels (comités, validations hiérarchiques longues) sont inopérants. L'entreprise doit adopter un mode de fonctionnement "commando" : centralisation de la décision et exécution rapide.

% ------------------------------------------------------------------------------
\section{Jargon de la Gestion de Crise}
% ------------------------------------------------------------------------------

\subsection{PCA vs PGC}

Il ne faut pas les confondre :

\begin{itemize}
    \item \textbf{PCA (Plan de Continuité d'Activité) :} "Comment on continue à travailler ?" (Focus : Opérations).
    \item \textbf{PGC (Plan de Gestion de Crise) :} "Comment on survit à la tempête ?" (Focus : Décision, Communication, Survie).
\end{itemize}

Le PCA est un outil technique ; le PGC est un outil de gouvernance.

\subsection{Zone de Turbulence}

Période d'intense agitation où l'information est brouillée, contradictoire et où les rumeurs fleurissent. C'est la phase la plus dangereuse.

% ------------------------------------------------------------------------------
\section{Cas LiZbiZ : Scénario "Catastrophe"}
% ------------------------------------------------------------------------------

\subsection{Mise en Situation}

\begin{boxlizbiz}
\textbf{Lundi matin, 08h30.}
Le titre boursier de LiZbiZ s'effondre de 15\% à l'ouverture.

\textbf{La cause :} Un tweet viral d'un hacker célèbre (@CyberGhost) publie un lien vers 10 000 dossiers médicaux de clients LiZbiZ (fuite de données santé).

\textbf{Réactions immédiates :}
\begin{itemize}
    \item Le standard téléphonique est saturé par les appels de clients paniqués.
    \item La boîte mail de la direction reçoit des menaces de recours collectifs.
    \item La DSI est introuvable (en réunion technique sur le Ransomware).
\end{itemize}

\textbf{Question :} Le PDG est au pied du mur. S'il ne réagit pas dans l'heure, la confiance sera définitivement perdue. C'est l'activation du PGC.
\end{boxlizbiz}

\subsection{Travaux Pratiques : Identification des Enjeux}

En groupes, les étudiants doivent lister les 3 priorités absolues pour le PDG dans les 30 prochaines minutes.

\textbf{Pistes de réflexion :}
\begin{enumerate}
    \item \textbf{Protéger les personnes :} Informer les clients concernés (Obligation RGPD).
    \item \textbf{Maîtriser l'information :} Ne pas répondre n'importe quoi aux journalistes. Préparer un "Holding Statement" (Déclaration de maintien).
    \item \textbf{Organiser le commandement :} Réunir la Cellule de Crise immédiatement.
\end{enumerate}

% ==============================================================================
% Fichier : 05_module_crise.tex (Extrait Chapitre 2)
% Description : Organisation de la Cellule de Crise
% ==============================================================================